%%%%%%%%%%%%%%%%%%%%%%%%%%%%%%%%%
%
%       EXERCÍCIO
%
%%%%%%%%%%%%%%%%%%%%%%%%%%%%%%%%%

\ifdefstring{\atividade}{prova}{%
    \renewcommand{\valorquestao}{\ValorQ\ pontos}
}%

\begin{exercicioBanco}[\valorquestao]
Suponha que o tempo (em horas) de falha de ventiladores em um computador pessoal possa ser modelado por uma distribuição exponencial, com média de 3300.
%
\begin{enumerate}[(a)]
    \item Qual a proporção de ventiladores que durará no mínimo 10.000 horas?
    %
    \item Qual a proporção de ventiladores que durará no máximo 7.000 horas?
    %
    \item Calcule \(P(X > 10000 \mid X > 8000)\) e \(P (X > 2000)\). Observe que esses resultados são iguais, consequência de uma importante propriedade da distribuição Exponencial chamada de Propriedade da Falta de Memória.
\end{enumerate}
%
\end{exercicioBanco}

